\chapter{Conclusion}\label{chap:conclusion}

\section{Summary of Contributions}
This dissertation presents the development and mechanical characterization of a novel composite scaffold system for meniscal tissue engineering, referred to as the. The key contributions of this work include:

\begin{itemize}
    \item \textbf{Gelatin Hydrogel Characterization:} Established a hydrogel system with tunable crosslinking and physiologically relevant water content and viscoelastic behavior, mimicking the time-dependent stress dissipation observed in native meniscus tissue.
    
    \item \textbf{Custom Mechanical Testing Apparatus:} Designed and validated a confined compression testing chamber, improving precision for low-modulus biomaterial testing while enabling real-time visualization and reproducibility.
    
    \item \textbf{Composite Scaffold Fabrication:} Integrated polycaprolactone (PCL) reinforcement into gelatin hydrogels using SLA-molded configurations. Patterns included aligned, circumferential, and 3D thermoformed PCL grids that enhanced compressive strength and anisotropic stress distribution.
    
    % \item \textbf{Mechanical Pattern Optimization:} Demonstrated the influence of strain rate and scaffold architecture on mechanical response, introducing the use of $E_{\text{confined}}$ to assess rate-dependent stiffness from 5–25\% strain.
    
    \item \textbf{Scalable, Cost-Effective Methodology:} Employed accessible additive manufacturing technologies (SLA and FDM) to enable personalized and structurally tuned scaffolds without reliance on expensive bioprinting equipment or complex cell-laden systems.
\end{itemize}

\section{Clinical Translation Potential}
This platform demonstrates significant potential for clinical translation in orthopedic tissue engineering:

\begin{itemize}
    \item \textbf{Customizability:} Scaffold architectures and mechanical properties can be tailored to patient-specific meniscal geometries and biomechanical demands using imaging data and reverse engineering workflows.
    
    \item \textbf{Biocompatibility and Tunability:} The hydrogel matrix provides a cell-supportive environment while allowing adjustment of crosslinking density, reinforcing the mechanical fidelity needed for implantation in weight-bearing joints.
    
    \item \textbf{Manufacturability:} The fabrication strategy avoids the scalability and cost barriers associated with bioprinting, offering a feasible solution for personalized medicine.
    
    \item \textbf{Scaffold Maturity:} Integration of aligned and circumferential reinforcement mimics the layered collagen organization of the native meniscus, improving load distribution and stress transformation during activity.
\end{itemize}

These findings represent a substantial step toward the development of biofunctional implants that replicate the mechanics and architecture of the human meniscus while remaining practical for clinical implementation.

\section{Closing Remarks on HANI and Tissue-Engineered Meniscus Design}


The modularity of the design, coupled with its adaptability to patient-specific anatomy and loading profiles, aligns with the growing demand for precision-medicine solutions in orthopedics. Future iterations may incorporate biological cues, controlled degradation profiles, and even in situ printing for surgical deployment.

Ultimately, this work reinforces the promise of engineered composite scaffolds to replicate the complexity of native fibrocartilage while providing practical pathways toward clinical use. It establishes not only the feasibility but the strategic framework for integrating multifunctional, personalized implants into the evolving landscape of regenerative medicine.
